\chapter{总结与展望}

\section{主要工作及结论}

本文面向微型无人机续航能力受限、复杂环境中长期驻留能力不足以及自主栖息可靠性有限等问题，围绕基于柔性附着结构的微型无人机系统设计与栖息控制展开研究，构建了一套融合轻量化飞行平台、视觉感知、飞行控制与柔性黏附执行机构的自主栖息系统，并通过仿真与实验对系统有效性进行了验证。本文的主要工作与研究结论如下：

\begin{enumerate}

\item 本文围绕微型无人机自主栖息任务，完成了无人机平台与柔性附着系统的整体设计与集成。
VSERF器件的制造工艺目前尚不成熟，黏附力有限。为此，本文在器件黏附需求与无人机动力之间寻求平衡，专门设计了一套适配方案——结构刚性有保障，动力学特性良好，能够可靠完成栖息任务。实验结果表明，该系统达到了预期设计指标，任务完成情况良好。与现有平台相比，本文方案在质量、惯量等参数上更为合理，整体结构更加紧凑，载重能力也略有提升，为后续进一步拓展应用留下了充足余量。
% 针对微型无人机在质量、推重比以及附着能力之间存在的强耦合约束，本文构建了一套基于 VSERF 的微型无人机栖息平台。在结构设计方面，通过轻量化碳纤维机架设计与底座拓扑优化，实现了无人机整体质量与结构刚度之间的平衡；同时结合推力测试与有限元分析，验证了平台在满足结构强度要求的前提下具有较高推力裕度，为自主栖息过程中姿态调整与接触稳定性提供了动力学基础。

% 此外，针对无人机顶部视觉感知需求与 VSERF 模块安装带来的偏心载荷问题，本文通过对称质量分布设计降低了系统质心偏移与重力诱导附加力矩，从结构层面提升了姿态控制稳定性。实验结果表明，所设计平台在满足轻量化要求的同时具备良好的飞行性能与系统集成能力。

\item 本文搭建了一套基于 AprilTag 的视觉伺服感知与控制框架，用于无人机自主栖息任务中的相对位姿估计与闭环控制。现实部署场景中，动作捕捉系统造价高昂、布置繁琐，难以实际应用。本文为此做了简化：地面站仅需一台普通笔记本电脑，定位标签用普通打印机打印即可，无需任何专用硬件。无人机通过实时识别 AprilTag，估计自身相对目标的位置与姿态，驱动控制器形成完整闭环，最终自主完成定位与栖息对准。与现有方法相比，本文方案在传感器配置极为简单的条件下，仍可实现亚分米级的定位精度，具备较强的实用价值。
% 考虑到微型无人机平台在机载算力与传感能力上的限制，本文采用地面站辅助计算架构，将视觉感知与导航决策部署于地面端，通过机载摄像头、图像回传链路以及 OpenCV 图像处理算法，实现基于 AprilTag 的目标识别与相对位姿估计。

% 在控制系统方面，本文构建了以姿态控制为核心、视觉反馈参与闭环调节的控制架构，并采用双层串级 PID 控制策略实现无人机姿态稳定与目标接近控制。针对栖息任务中存在的视觉目标局部遮挡与短时丢失问题，对视觉可观测性边界与安全工作区域进行了分析。研究结果表明，合理设计视觉目标尺度及控制策略后，系统能够在视觉信息退化情况下维持短时稳定飞行，为最终接触与黏附动作提供可靠控制基础。

\item 基于 MuJoCo 仿真平台构建了无人机自主栖息验证环境，并完成系统级可行性分析。
为降低实机实验风险并验证控制方案有效性，本文基于 MuJoCo 物理引擎建立了无人机栖息任务仿真环境，集成无人机动力学模型、视觉反馈模块、控制器模块以及黏附执行机构模型，实现了从起飞、目标识别、视觉接近到最终附着的完整任务流程验证。仿真结果表明，在一定程度的视觉信息退化或短时视野缺失情况下，系统仍能够依赖姿态闭环控制与动力学惯性保持有限稳定性，并最终完成预期接近与附着任务。这说明本文所构建控制框架在复杂感知条件下具备一定鲁棒性，同时验证了视觉伺服与飞行控制融合策略的工程可行性。

\item 完成了 VSERF 柔性附着执行机构与无人机飞行控制系统的协同集成，并验证了系统自主栖息能力。
本文完成了 VSERF 柔性黏附器件、高压驱动电路、无人机本体的系统集成，通过光耦隔离实现低压飞控系统对高压驱动回路的安全控制，并结合飞控底层 GPIO 功能重映射实现附着动作触发控制。实验结果表明，VSERF 器件能够在满足质量与能耗约束条件下提供足够附着能力，使无人机在目标表面实现稳定附着。结合自主飞行控制与视觉伺服系统，所构建无人机平台能够完成目标接近、姿态调整与附着动作执行，起飞后不再依赖人工操作，验证了本文提出的柔性附着无人机系统在自主栖息任务中的可行性。

\end{enumerate}

\vspace{5pt}

综上所述，本文围绕微型无人机自主栖息问题，从系统结构设计、柔性附着执行机构、视觉感知与控制策略等方面开展了系统研究，构建了一套具备自主接近与附着能力的微型无人机栖息系统。研究结果表明，通过轻量化结构设计、视觉伺服控制与 VSERF 柔性附着机制的协同融合，可以在有限载荷与受限感知条件下实现无人机自主栖息任务，为提升微型无人机续航能力与长期驻留能力提供了一种具有工程应用潜力的技术路径。

\section{研究展望}

尽管本文完成了基于 VSERF 柔性附着机构的微型无人机自主栖息系统设计，并验证了其在特定实验条件下的可行性，但仍存在进一步优化与拓展空间。未来研究可重点围绕以下方向展开：

\begin{enumerate}

\item 当前实验主要针对目标平面法向与重力方向一致的典型场景开展，即无人机在 AprilTag 引导下完成目标下方悬停、姿态调整与附着过程。为进一步提升系统在复杂环境中的适应能力，后续研究可扩展至倾斜平面甚至任意姿态目标表面的栖息任务，不再依赖AprilTag的引导，重点研究不同表面法向条件下的视觉感知、姿态控制与附着稳定性问题，并分析重力分量变化对接触动力学及附着性能的影响，以增强系统在非结构化环境中的应用能力。

\item 本文所实现的自主接近与悬停控制主要面向静止目标，而所构建的控制框架本质上具备对目标相对运动进行补偿的能力。未来研究可进一步考虑目标平面存在一定平动速度的动态场景，探索无人机在地速方向非零条件下的自主起飞、接近与附着控制问题，实现对移动目标或运动平台的稳定栖息与脱附，从而拓展系统在移动载体协同、自主巡检及动态环境作业等任务中的应用潜力。

\item 当前系统受限于微型无人机平台的载荷能力与柔性附着器件性能边界，在负载能力方面仍存在一定约束。未来研究可从提升 VSERF 器件附着性能与系统集成能力两个方向展开：一方面，可通过优化材料配方、结构设计或驱动方式，提高单个器件的附着强度与稳定性；另一方面，可探索多个附着器件的协同集成方案，以提升系统整体承载能力，从而支持更复杂功能模块的搭载，并进一步扩展无人机在复杂任务场景中的应用能力。

\end{enumerate}
